\documentclass[11pt]{article}
\usepackage[french]{babel}
\frenchbsetup{og=«,fg=»}
\usepackage{geometry}                % See geometry.pdf to learn the layout options. There are lots.
\geometry{a4paper}
\usepackage{graphicx}
\usepackage{pifont}
\usepackage{enumitem}
\usepackage{xspace}
\usepackage{tocloft}
\usepackage{mathtools}
\usepackage{mathdots}
\usepackage{amssymb}
\usepackage{amsmath}
\usepackage{amsfonts}
\usepackage{cancel}
\usepackage{hyperref}
\hypersetup{
    colorlinks=true,
    linkcolor=blue,
    filecolor=magenta,      
    urlcolor=cyan,
    pdftitle={Overleaf Example},
    pdfpagemode=FullScreen,
    }

\DeclareGraphicsRule{.tif}{png}{.png}{`convert #1 `dirname #1`/`basename #1 .tif`.png}
\usepackage[ddmmyyyy]{datetime}
\usepackage{rotating}
\usepackage{lscape}
\usepackage{relsize}
\usepackage{longtable}
\usepackage{array}
\usepackage{xcolor}
\usepackage{listings}
\usepackage{listingsutf8}
\definecolor{bgcollst1}{rgb}{0.95,0.95,0.95}
\definecolor{bgcollst2}{rgb}{0.7,0.7,0.7}
\lstset{inputencoding=utf8/latin1}
\lstset{literate={œ}{{\oe}}1}
%\lstset{numbers=left, numberstyle=\tiny, stepnumber=1, numbersep=5pt}
\lstset{
%	inputencoding=latin1,
  language=C++,
  flexiblecolumns=true,
  basicstyle=\ttfamily\footnotesize,
  keywordstyle=\color{blue}\bfseries,
  commentstyle=\color{green!60!black},
  stringstyle=\color{red},
  numbers=left,
  numberstyle=\tiny\color{gray},
  breaklines=true,
  frame=single,
  showstringspaces=false,
  tabsize=2,
  backgroundcolor=\color{bgcollst1},
  frame=shadowbox,
  framerule=0.5pt,
  framesep=0.5pt,
  rulesep=1pt,
  framexleftmargin=1mm,
  framexrightmargin=1mm,
  framextopmargin=1mm,
  framexbottommargin=1mm,
  fillcolor=\color{black},
  rulecolor=\color{black},
  rulesepcolor=\color{bgcollst2}
}
%\lstset{
%	inputencoding=latin1,
%	language=C++,
%	flexiblecolumns=true,
%	morekeywords={append, as, extend, insert, sort, remove, reverse, keys, values, items, write, close, read, with, self, cls, time copy, dtype, yield, match, case},
%	basicstyle=\scriptsize\ttfamily,
%	showstringspaces=false,
%	framexleftmargin=1mm,
%	framexrightmargin=1mm,
%	framextopmargin=1mm,
%	framexbottommargin=1mm,
%	tabsize=2,
%	backgroundcolor=\color{bgcollst1},
%	frame=shadowbox,
%	framerule=0.5pt,
%	framesep=0.5pt,
%	rulesep=1pt,
%	fillcolor=\color{black},
%	rulecolor=\color{black},
%	rulesepcolor=\color{bgcollst2},
%	keywordstyle=\color{black}\bfseries\underbar,
%}
 \lstset{%
	inputencoding=utf8,
	extendedchars=true,
	literate=%
		{é}{{\'e}}{1}%
		{è}{{\`e}}{1}%
		{à}{{\`a }}{1}%
		{ç}{{\c{c}}}{1}%
		{œ}{{\oe}}{1}%
		{ù}{{\`u}}{1}%
		{É}{{\'E}}{1}%
		{È}{{\`E}}{1}%
		{À}{{\`A}}{1}%
		{Ç}{{\c{C}}}{1}%
		{Œ}{{\OE}}{1}%
		{Ê}{{\^E}}{1}%
		{ê}{{\^e}}{1}%
		{î}{{\^i}}{1}%
		{ï}{{\"{\i}}}{1}%
		{ô}{{\^o}}{1}%
		{û}{{\^u}}{1}%
		{ë}{{\¨{e}}}1
		{û}{{\^{u}}}1
		{â}{{\^{a}}}1
		{Â}{{\^{A}}}1
		{Î}{{\^{I}}}1
		{Ï}{{\"{\I}}}{1}
		{·}{{\textperiodcentered}}{1}	
		{²}{{\texttwosuperior}}{1}
		{³}{{\textthreesuperior}}{1}
		{≤}{{$\leqslant\ $}}{1}
		{≥}{{$\geqslant\ $}}{1}
		{∈}{{$\in\ $}}{1}
		{π}{{$\pi \ $}}{1}
		{λ}{{$\lambda \ $}}{1}
%		{ℝ}{{$\lambda \ $}}{1}
}
\newcommand{\N}{\mathbf{\mathds{N}}}
\newcommand{\Z}{\mathbf{\mathds{Z}}}
\newcommand{\Q}{\mathbf{\mathds{Q}}}
\newcommand{\R}{\mathbf{\mathds{R}}}
\newcommand{\COMPLEX}{\mathbf{\mathds{C}}}
\newcommand{\K}{\mathbf{\mathds{K}}}

\newcommand{\ds}{\displaystyle}
\newcommand{\mb}[1]{\mathbf{#1}}
\newcommand{\vb}[1]{\mathbf{#1}}

\DeclareMathOperator*{\mmod}{\mathrm{\mathbf{mod}}}
%\DeclareMathOperator*{\aabs}{\mathrm{abs}}
\DeclareMathOperator*{\kker}{\mathrm{\mathbf{ker}}}
\DeclareMathOperator*{\iimg}{\mathrm{\mathbf{img}}}
\DeclareMathOperator*{\ddim}{\mathrm{\mathbf{dim}}}
\DeclareMathOperator*{\iinf}{\mathrm{\mathbf{inf}}}
\DeclareMathOperator*{\ssup}{\mathrm{\mathbf{sup}}}
\DeclareMathOperator*{\mmin}{\mathrm{\mathbf{min}}}
\DeclareMathOperator*{\mmax}{\mathrm{\mathbf{max}}}
\DeclareMathOperator*{\ddeg}{\mathrm{\mathbf{deg}}}
\DeclareMathOperator*{\vval}{\mathrm{\mathbf{val}}}
\DeclareMathOperator*{\rres}{\mathrm{\mathbf{res}}}
\DeclareMathOperator*{\ddis}{\mathrm{\mathbf{disc}}}
\DeclareMathOperator*{\pppcm}{\mathrm{\mathbf{ppcm}}}
\DeclareMathOperator*{\ppgcd}{\mathrm{\mathbf{pgcd}}}
\DeclareMathOperator*{\ddet}{\mathrm{\mathbf{det}}}
\DeclareMathOperator*{\aadj}{\mathrm{\mathbf{adj}}}
\DeclareMathOperator*{\iinv}{\mathrm{\mathbf{inv}}}
\DeclareMathOperator*{\ttra}{\mathrm{\mathbf{tr}}}
\DeclareMathOperator*{\rrng}{\mathrm{\mathbf{rang}}}
\DeclareMathOperator*{\hhes}{\mathrm{\mathbf{hess}}}
\DeclareMathOperator*{\jjac}{\mathrm{\mathbf{jac}}}
\DeclareMathOperator*{\ggra}{\mathrm{\mathbf{grad}}}
\DeclareMathOperator*{\ddiv}{\mathrm{\mathbf{div}}}
\DeclareMathOperator*{\rrot}{\mathrm{\mathbf{rot}}}
\DeclareMathOperator*{\dist}{\mathrm{\mathbf{dist}}}
\DeclareMathOperator*{\ccov}{\mathrm{\mathbf{cov}}}
\DeclareMathOperator*{\vvar}{\mathrm{\mathbf{var}}}

\let\mathds\mathbb
\linespread{1.2}
\renewcommand{\arraystretch}{1.5}

\interfootnotelinepenalty=5000
\widowpenalty=10000
\clubpenalty=10000
\hyphenpenalty=10000
\uchyph=0
\thinmuskip=1mu
\delimiterfactor=950
\raggedbottom
\tolerance=1
\emergencystretch=\maxdimen
\setlength\parindent{0pt}

\title{\Huge Friandises \textcolor{red}{\textbf{MATRICIELLES}} en \textbf{C++}}

\author{Peter Philippe}

\begin{document}
\maketitle
Ces quelques pages proposent diff\'erents calculs num\'eriques (d\'eterminant, norme, inverse, polyn\^ome caract\'eristique) et des d\'ecompositions (\textbf{Cholesky}, \textbf{LDLT}, \textbf{LU},
\textbf{QR}, \textbf{Diagonalisation}, \textbf{Schur}, \textbf{Jordan-Chevalley}, \textbf{Iwasawa} et \textbf{Polaire}),
pour des matrices carr\'ees d'ordre~$2$ \`a coefficients dans~$\R$.\\

L'ordre \'etant tr\`es petit, les calculs des diff\'erents algorithmes propos\'es seront donc scandaleusement simplifi\'es dans le code source. La notation utilis\'ee pour les coefficients de
la matrice est classique~:
\[
\mb{A}=
\begin{pmatrix}
a_{1,1} & a_{1,2} \\
a_{2,1} & a_{2,2}
\end{pmatrix}
\]
Parmi les diff\'erentes m\'ethodes exhib\'ees, la d\'ecomposition de Jordan peut surprendre par son absence, mais avec un ordre~$2$, il n'a pas \'et\'e utile
de la mentionner, car cela revient \`a diagonaliser une matrice. De plus, avec cet ordre \og scolaire \fg, la v\'erification permettant de qualifier qu'une matrice
est structur\'ee n'a pas \'et\'e ajout\'ee. Ici, le produit matriciel de \texttt{operator * (...)} est celui de Strassen (d'autres sont possibles comme celui de Winograd
ou de Pan).\\

\tiny{\textcolor{red}{\textit{(Mieux vaut avoir des bases en alg\`ebre lin\'eaire pour \^etre \`a l'aise...)}}}\normalsize

\section{Définitions succintes des méthodes de la classe \texttt{Matri2x2}}

\begin{description}[leftmargin=4cm, style=nextline]

\item[\texttt{trace()}] 
La trace est un invariant de similitude par transposition (changement de base), il s'agit simplement d'un scalaire calcul\'e en sommant tous les coefficients
de la diagonale de la matrice~$\ttra(\mb{A}) = \sum_{i}a_{i,i}$.\\

\item[\texttt{determinant()}] 
Le d\'eterminant est une fonction antisym\'etrique (ici il se r\'eduit \`a $|\mb{A}|=a_{1,1}a_{2,2}-a_{2,1}a_{1,2}$), lorsqu'il est non nul, l'inversibilit\'e de la matrice sera assur\'ee,
elle sera alors qualifi\'ee de \og r\'eguli\`ere \fg, donc de rang plein, dans le cas contraire, elle sera dite \og singuli\`ere \fg. Il intervient notamment dans le calcul du polynôme caractéristique.\\

\item[\texttt{norm()}] 
La norme de Frobenius, encore appel\'ee norme euclidienne, est une fonction retournant une valeur positive, cela permet de mesurer la longueur de la matrice dans
son espace et ainsi de pouvoir en comparer la similitude avec une autre matrice, la formule est~$\|\mb{A}\|_{2} = \sqrt{\sum_{i}\sum_{j} |a_{i,j}|^{2}}$. Cette mesure
est indispensable en analyse num\'erique matricielle.\\

\item[\texttt{is\_singular()}] 
Une matrice singulière n'est pas inversible, il y a alors au moins une ligne ou une colonne proportionnelle \`a une autre (famille de vecteurs li\'es), son noyau est non trivial.
Le rang de la matrice est alors strictement inf\'erieur \`a son ordre.\\

\item[\texttt{is\_symmetric()}] 
Une matrice $\mb{A}$ est de type sym\'etrique si $\mb{A} = \mb{A}^{\textsc{T}}$, ses \'el\'ements sont sym\'etriques par rapport \`a la diagonale $(a_{ij} = a_{ji})$, ici cela
se r\'esume \`a~$b=c$. Ces matrices sont toujours diagonalisables et toutes leurs valeurs propres sont r\'eelles.\\

\item[\texttt{is\_nilpotent()}] 
Une matrice $\mb{A}$ est dite nilpotente si $\mb{A}^{k-1} \neq 0$ et $\mb{A}^k = 0$ pour un certain $k \in \mathbb{N}^*$; pour un ordre~$2$, cela \'equivaut donc
\`a~$\mb{A}^2 = 0$, sa trace et son d\'eterminant sont nuls.\\

\item[\texttt{is\_idempotent()}] 
Une matrice $\mb{A}$ idempotente est class\'ee comme telle lorsque $\mb{A}^n = \mb{A}$ ($n$ correspond \`a l'ordre de $\mb{A}$, donc ici $\mb{A}^2 = \mb{A}$),
ce qui correspond \`a une projection lin\'eaire; ses valeurs propres sont $0$ ou $1$.\\

\item[\texttt{is\_involutory()}] 
On dit d'une matrice carr\'ee $\mb{A}$ qu'elle est involutive lorsque~$\mb{A}\mb{A} = \mb{I}$, son d\'eterminant vaut~$|\mb{A}|=1$ ou~$|\mb{A}|=-1$, elle est alors sa propre inverse.\\

\item[\texttt{is\_orthogonal()}] 
Une matrice $\mb{A}$ est orthogonale si $\mb{A}\mb{A}^{\textsc{T}}=\mb{A}\mb{A}^{-1} = \mb{I}$ (son inverse est donc \'egal \`a sa transpos\'ee, ce qui est nettement plus rapide \`a calculer !).
Ses colonnes forment une base orthonorm\'ee. Par exemple, la matrice de rotation~$\mb{A}=\begin{psmallmatrix}\cos\,\theta&-\sin\,\theta \\ \sin\,\theta&\cos\,\theta\end{psmallmatrix}$
dans~$\R^{2}$ est un \'el\'ement du groupe special orthogonal not\'e~$\textsc{SL}_{2}(\R)$.\\

\item[\texttt{is\_upper\_unipotent()}] 
Vérifie que la matrice est triangulaire supérieure avec des $1$ sur la diagonale ($a_{1,1} =a_{2,2} =1$, $a_{2,1} =0$).\\

\item[\texttt{is\_diagonal\_positive()}] 
Teste si la matrice est diagonale ($a_{1,2} =a_{2,1} =0$) et si ses deux coefficients diagonaux sont strictement positifs.\\

\item[\texttt{transpose()}] 
Retourne la matrice transpos\'ee $\mb{A}^{\textsc{T}}$ telle que $a_{i,j} \to a_{j,i}$.\\

\item[\texttt{adjugate()}] 
Retourne la matrice adjointe $\mb{A}^{*}$, utilisée pour le calcul de l'inverse~$\mb{A}^{-1}=\ddet(\mb{A})^{-1}\mb{A}^{*}$.\\

\item[\texttt{inverse()}] 
Une matrice carr\'ee~$\mb{A}$ d'ordre~$n$ \`a coefficients dans~$\K$ est inversible si et seulement si son d\'eterminant est non nul~$|\mb{A}|\neq0$,
son rang~$k$ est alors n\'ecessairement \'egal \`a son ordre~$n$. L'unique solution de~$\mb{A}\vb{x}=\vb{0}$ est le vecteur~$\vb{x}=\vb{0}$,
cela sous-entend que la famille form\'ee par tous ses vecteurs colonnes est lin\'eairement ind\'ependante (famille libre). L'ensemble des matrices~$\mb{A}$ inversibles
d'ordre~$n$ (et diagonalisables) \`a coefficients dans~$\R$ est un groupe lin\'eaire not\'e~$\textsc{GL}_{n}(\R)$. Dans le cas o\`u la condition de d\'eterminant
non nul (ou de rang plein) est satisfaite, la matrice~$\mb{A}$ d'ordre~$n$ poss\`ede une matrice inverse, unique, telle que~$\mb{A}\mb{A}^{-1} = \mb{I}$.\\

\item[\texttt{Cholesky()}] 
Décompose une matrice symétrique et définie positive (SDP) en $\mb{A} = \mb{LL}^{\textsc{T}}$, la fonction retourne l'unique matrice triangulaire inférieure~$\mb{L}$ \`a \'el\'ements diagonaux positifs.\\

\item[\texttt{LDLT()}] 
Lorsque la m\'ethode de Cholesky \'echoue du fait que la matrice n'est pas (SDP), il faut utiliser la m\'ethode $\mb{A} = \mb{L}\mb{D}\mb{L}^{\textsc{T}}$ avec~$\mb{L}$
triangulaire inf\'erieure unitaire et~$\mb{D}$ diagonale, car elle n'utilise plus les racines carr\'ees (ce qui la rend aussi plus stable numériquement). Ici aussi, la matrice~$\mb{L}$ est unique.\\

\item[\texttt{LU()}] 
Cette factorisation $\mb{A} = \mb{LU}$ ne s'applique que si $\mb{A}$ est inversible. Cela factorise la matrice en une triangulaire inf\'erieure~$\mb{L}$ (avec une diagonale unitaire)
et une triangulaire sup\'erieure $\mb{U}$. Cet algorithme permet de r\'esoudre des syst\`emes lin\'eaires du type~$\mb{L}\mb{U}\vb{x}=\vb{b}$ en recherchant~$\vb{x} = \mb{U}^{-1}(\mb{L}^{-1}\vb{b})$.
Cette m\'ethode permet aussi le calcul du d\'eterminant car $|\mb{A}|=|\mb{U}|$, mais aussi celui de l'inverse de cette fa\c{c}on~:
$$\mb{A}^{-1} = \mb{U}^{-1}\mb{L}^{-1}$$
Si le premier pivot $a_{1,1}$ est nul, il faut alors proc\'eder \`a une permutation des lignes.\\

\item[\texttt{QR()}] 
La d\'ecomposition unique $\mb{A} = \mb{Q}\mb{R}$ avec $\mb{Q}$ orthogonale et unitaire (ce qui sous-entend ici diagonalisable) et~$\mb{R}$ triangulaire sup\'erieure telles
que~$\mb{Q}\mb{Q}^{\textsc{T}}=\mb{Q}\mb{Q}^{-1}=\mb{I}$ (donc~$\mb{Q}^{\textsc{T}}=\mb{Q}^{-1}$). Cette d\'ecomposition peut \^etre obtenue par diff\'erents proc\'ed\'es
(ici, c'est celui de Gram-Schmidt modifi\'e), il s'agit d'un algorithme fondamental en alg\`ebre lin\'eaire num\'erique. Un cas bien classique de matrice orthogonale est la matrice
(non homog\`ene)~$\mb{R}(\theta)$ du groupe~$\text{SO}(2)$ (groupe \underline{S}p\'ecial \underline{O}rthogonal des matrices d'ordre 2) utilis\'ee pour une rotation
d'angle~$\theta\in[0,2\pi[$ (lorsque~$\theta=0$ ou~$\theta=\pi$, il y a~$\mb{R}(\theta)=\mb{I}$) dans le plan euclidien~$\R^{2}$.\\

\item[\texttt{characteristic\_polynomial()}] 
Retourne le polynôme caract\'eristique unitaire (dont le degr\'e~$n$ correspond \`a l'ordre de~$\mb{A}$) tel que~:\\$\mathcal{P}_{\mb{A}}(\lambda) = |\mb{A}-\lambda\mb{I}|= (-1)^{n}\lambda^2 - \ttra(\mb{A})\lambda + \ddet(\mb{A})$.
La diagonalisation ou encore le th\'eor\`eme de Cayley-Hamilton utilisent le polynôme caract\'eristique.\\

\item[\texttt{diagonalization()}] 
Si la matrice est diagonalisable, avec~$\mathcal{P}_{\mb{A}}$ scind\'e et \`a racines simples distinctes dans~$\COMPLEX$, la m\'ethode retourne~$(\mb{P}, \mb{D}, \mb{P}^{-1})$ où~$\mb{D}$ est diagonale (compos\'ee des valeurs propres~$\lambda$
de~$\mb{A}$) et~$\mb{P}$ la matrice des vecteurs propres (colonnes) de~$\mb{A}$.\\

\item[\texttt{schur()}]
Cette d\'ecomposition $\mb{A} = \mb{S}\mb{D}\mb{S}^{\textsc{T}}$ avec $\mb{S}$ unitaire (orthogonale pour les r\'eelles) et~$\mb{D}$ diagonale ou triangulaire sup\'erieure
(ou m\^eme \^etre une matrice de rotation dans le cas complexe); permet de r\'ev\'eler la structure spectrale de la matrice.
\[
\begin{split}
\mb{A} &=
\begin{pmatrix}
a_{1,1} & a_{1,2} \\
a_{2,1} & a_{2,2}
\end{pmatrix} \\
&=
\begin{pmatrix*}[r]
1 & 0 \\
a_{1,1}^{-1}a_{2,1} & 1
\end{pmatrix*}
\begin{pmatrix*}[r]
a_{1,1} & 0 \\
0 & \mb{S}
\end{pmatrix*}
\begin{pmatrix*}[r]
1 & a_{1,1}a_{1,2}^{-1} \\
0 & 1
\end{pmatrix*}
\end{split}
\]
avec $\mb{S}=a_{2,2} - a_{2,1}a_{1,1}^{-1}a_{1,2}$.\\

\item[\texttt{jordan\_chevalley()}] 
Dans un $\K$-espace vectoriel de dimension finie, pour certaines matrices carrées $\mb{A} \in \mathcal{M}_n(\mathbb{C})$ de polyn\^ome minimal\footnote{Attention, $\mathcal{P}$ est scind\'e (produits de polyn\^omes de degr\'e~1) si et seulement si le polyn\^ome caract\'eristique l'est \'egalement.}~$\mathcal{P}(\mb{A})$ scind\'e,
il existe une unique paire de matrices $\mb{D}$ (diagonalisable sur $\mathbb{C}$ uniquement si $\mathcal{P}(\mb{A})$ est scind\'e) et $\mb{N}$ (nilpotente) telles que :
$$\mb{A} = \mb{D} + \mb{N} \qquad \mb{D}\mb{N}=\mb{N}\mb{D} \qquad \mathcal{P}(\mb{A})=\mathcal{P}(\mb{D})$$
Les matrices d'ordre 2 non diagonalisables v\'erifient $\mb{N}=\mb{A}-\mb{D}$ et~$\mb{D}=\lambda\mb{I}$ ($\lambda$ est une valeur propre de~$\mb{A}$)
ainsi que~$\mb{N}=\mb{A}-\lambda\mb{I}\).\\

\item[\texttt{iwasawa()}] 
Pour des matrices carr\'ees d'ordre 2, cette d\'ecomposition d'une matrice $\mb{A}$ prend la forme suivante~:
$$\mb{A} = \mb{Q}\mb{D}\mb{N}$$
avec $\mb{Q}$ une matrice orthogonale, $\mb{D}$ une matrice diagonale dont les coefficients sont strictement positifs et~$\mb{N}$ une matrice unipotente unitaire sup\'erieure (donc triangulaire).
La matrice $\mb{Q}$ s'obtient par le proc\'ed\'e de Gram-Schmidt, ses deux vecteurs colonnes $q_1$ et $q_2$ sont~:
$$q_1 = \dfrac{(a_{11},a_{21})^{\textsc{T}}}{\sqrt{a_{11}^2+a_{21}^2}} \text{ et } q_2 = \dfrac{(a_{12},a_{22})^{\textsc{T}}-(q_1^{\textsc{T}}(a_{12},a_{22})^{\textsc{T}})q_1}{\|(a_{12},a_{22})^{\textsc{T}}-(q_1^{\textsc{T}}(a_{12},a_{22})^{\textsc{T}})q_1\|}$$
La diagonale de $\mb{D}$ est compos\'ee avec les \'el\'ements de la diagonale de la matrice~$\mb{T}$~:
$$\mb{T} = \mb{Q}^{\textsc{T}}\mb{A} = \begin{pmatrix}q_1^{\textsc{T}}(a_{11},a_{21})^{\textsc{T}} & q_2^{\textsc{T}}(a_{12},a_{22})^{\textsc{T}} \\ q_1^{\textsc{T}}(a_{11},a_{21})^{\textsc{T}} & q_2^{\textsc{T}}(a_{12},a_{22})^{\textsc{T}}\end{pmatrix}$$
C'est-\`a-dire~:
$$\mb{D} = \begin{pmatrix} t_1 & 0 \\ 0 & t_2\end{pmatrix}$$
Et $\mb{N}$ est construite comme ceci~:
$$\mb{N} = \begin{pmatrix} 1 &t_{1,2}/d_1 \\ 0 & 1\end{pmatrix}$$\\

\item[\texttt{polar\_decomposition()}] 
Une matrice carr\'ee $\mb{A}$ \`a coefficients dans~$\R$ peut \^etre d\'ecompos\'ee en un produit de matrices $\mb{A}=\mb{Q}\mb{S}$ avec~$\mb{Q}$ une matrice orthogonale et~$\mb{S}$ une matrice r\'eelle sym\'etrique positive.
Lorsque~$\mb{A}$ est inversible, les matrices~$\mb{S}$ et~$\mb{Q}$ sont uniques. De plus, comme~$\mb{A}=\mb{Q}\mb{S}$ alors~$\mb{A}^{\textsc{T}}=\mb{S}\mb{Q}^{\textsc{T}}$ et donc~$\mb{A}^{\textsc{T}}\mb{A} = \mb{S}^{2}$.
(Ne pas omettre que $\mb{A}^{\textsc{T}}\mb{A}\neq\mb{A}\mb{A}^{\textsc{T}}$.) Le calcul est $\mb{P} = \sqrt{\mb{A}^{\textsc{T}}\mb{A}}$ puis $\mb{Q} = \mb{A}\mb{S}^{-1}$.
La m\'ethode \texttt{polar\_decomposition\_general()} retourne $\mb{A} = \mb{Q}\mb{S}$ avec $\mb{Q}$ orthogonale et $\mb{S}$ symétrique positive.\\
\end{description}

\section{Impl\'ementation en C++}
Comme \`a mon habitude, ce code de presque 1300 lignes est bien \'evidemment perfectible (j'aurais quand m\^eme pu faire l'effort de rajouter le choix du type dans l'argument au lieu de le laisser
au tout d\'ebut de la fonction \texttt{main()}), il ne comporte qu'assez peu de commentaires, sa compilation avec GCC 15.2 s'effectue sans le moindre avertissement. L'utilisation est fort simple puisqu'il
ne faut saisir que les 4 coefficients de la matrice (r\'eels ou entiers) s\'epar\'es par une espace, si aucune valeur n'est saisie, une matrice \`a coefficients entiers sera g\'en\'er\'ee automatiquement.
\begin{lstlisting}
// g++-mp-15 -march=native -O3 -std=c++23 -lm -Wall -Wpedantic -Wextra -Werror -Wfatal-errors matri2x2.cpp -o matri2x2
#include <iostream>
#include <memory>
#include <complex>
#include <exception>
#include <stdexcept>
#include <limits>
#include <optional>
#include <utility>
#include <array>
#include <variant>
#include <tuple>
#include <string>
#include <cmath>
#include <random>
#include <print>
#include <format>
#include <concepts>
#include <string_view>
#include <valarray>

template<typename T>
struct std::formatter<std::complex<T>> {
    constexpr auto parse(std::format_parse_context& context) {
        return context.begin();
    }
    
    auto format(const std::complex<T>& c, std::format_context& context) const {
        if (c.imag() == T(0)) {
            return std::format_to(context.out(), "{}", c.real());
        } else if (c.real() == T(0)) {
            return std::format_to(context.out(), "{}i", c.imag());
        } else if (c.imag() < T(0)) {
            return std::format_to(context.out(), "{}{}i", c.real(), c.imag());
        } else {
            return std::format_to(context.out(), "{}+{}i", c.real(), c.imag());
        }
    }
};

template <typename T> concept MatrixType =
    std::is_arithmetic_v<T> ||
    std::is_same_v<T, std::complex<float>> ||
    std::is_same_v<T, std::complex<double>> ||
    std::is_same_v<T, std::complex<long double>>;

template <MatrixType T = float> 
class Matri2x2 final {
public:    
    using k_type = std::conditional_t<
        std::is_same_v<T, std::complex<float>>, float,
        std::conditional_t<std::is_same_v<T, std::complex<double>>, double,
        std::conditional_t<std::is_same_v<T, std::complex<long double>>, long double, T>>
    >;
    
    static constexpr k_type _eps = std::numeric_limits<k_type>::epsilon();
    static constexpr size_t N = 2;
    
    static constexpr bool isComplex = 
        std::is_same_v<T, std::complex<float>> ||
        std::is_same_v<T, std::complex<double>> ||
        std::is_same_v<T, std::complex<long double>>;

    using complexType =
        std::conditional_t<std::is_same_v<T, float>, std::complex<float>,
        std::conditional_t<std::is_same_v<T, double>, std::complex<double>,
        std::conditional_t<std::is_same_v<T, long double>,  std::complex<long double>, T>
        >>;

private:
    alignas(64) std::valarray<T> _data;
    mutable std::optional<bool> _singular, _symmetric, _nilpotent, _idempotent, _involutory, _orthogonal, _upper_unipotent, _diagonal_positive;

    constexpr T& a() { return _data[0]; }
    constexpr T& b() { return _data[1]; }
    constexpr T& c() { return _data[2]; }
    constexpr T& d() { return _data[3]; }
    
    constexpr const T& a() const { return _data[0]; }
    constexpr const T& b() const { return _data[1]; }
    constexpr const T& c() const { return _data[2]; }
    constexpr const T& d() const { return _data[3]; }

public:
    explicit constexpr Matri2x2() noexcept : _data(T(0), 4) {
        a() = T(1);
        d() = T(1);
    }
    
    explicit constexpr Matri2x2(const T& a_val, const T& b_val, const T& c_val, const T& d_val) noexcept 
        : _data(4) {
        a() = a_val;
        b() = b_val;
        c() = c_val;
        d() = d_val;
    }

    explicit constexpr Matri2x2(const std::array<T, 4>& stdarr) noexcept
    : _data(stdarr.data(), 4) {}

    explicit constexpr Matri2x2(std::array<T, 4>&& stdarr) noexcept
    : _data(std::move(stdarr)) {}

    template<MatrixType U>
    explicit constexpr Matri2x2(const Matri2x2<U>& mat) noexcept
        : _data{static_cast<T>(mat[0]), static_cast<T>(mat[1]), static_cast<T>(mat[2]), static_cast<T>(mat[3])} {}

    constexpr Matri2x2(const Matri2x2<T>&) noexcept = default;
    constexpr Matri2x2<T>& operator=(const Matri2x2<T>&) noexcept = default;

    constexpr Matri2x2(Matri2x2<T>&&) noexcept = default;
    constexpr Matri2x2<T>& operator=(Matri2x2<T>&&) noexcept = default;

    // ~Matri2x2() = default; -> inutile 
    
    constexpr T& operator () (size_t i, size_t j) {
        if (i >= N || j >= N)
            throw std::out_of_range("(i,j) out of matrix!");
        [[assume(i < N && j < N)]];
        return _data[i * N + j];
        // return (i==0 && j==0) ? a_ : (i==0 && j==1) ? b_ : (i==1 && j==0) ? c_ : d_;
    }
    
    constexpr T operator () (size_t i, size_t j) const {
        if (i >= N || j >= N)
        throw std::out_of_range("(i,j) out of matrix !");
        [[assume(i < N && j < N)]];
        return _data[i * N + j];
        // return (i==0 && j==0) ? a_ : (i==0 && j==1) ? b_ : (i==1 && j==0) ? c_ : d_;
    }

    constexpr const T& operator[](size_t i) const noexcept {
        return _data[i];
    }
    constexpr T& operator[](size_t i) noexcept {
        return _data[i];
    }

    [[nodiscard]] constexpr bool operator == (const Matri2x2<T>& m) const noexcept {
        if constexpr (std::is_integral_v<T>) {
            return a() == m.a() && b() == m.b() && c() == m.c() && d() == m.d();
        } else if constexpr (isComplex) {
            return std::abs(a() - m.a()) <= _eps && std::abs(b() - m.b()) <= _eps &&
                   std::abs(c() - m.c()) <= _eps && std::abs(d() - m.d()) <= _eps;
        } else {
            return std::abs(a() - m.a()) <= _eps && std::abs(b() - m.b()) <= _eps &&
                   std::abs(c() - m.c()) <= _eps && std::abs(d() - m.d()) <= _eps;
        }
    }

    [[nodiscard]] constexpr bool operator != (const Matri2x2<T>& m) const noexcept {
        return !(*this == m);
    }

    [[nodiscard]] constexpr Matri2x2<T> operator + (const Matri2x2<T>& m) const noexcept {
        Matri2x2<T> result;
        result._data = _data + m._data;
        return result;
    }

    [[nodiscard]] constexpr Matri2x2<T> operator - (const Matri2x2<T>& m) const noexcept {
        Matri2x2<T> result;
        result._data = _data - m._data;
        return result;
    }

    [[nodiscard]] constexpr Matri2x2<T> operator * (const Matri2x2<T>& m) const noexcept {
        if consteval {
            return Matri2x2<T>(
                a() * m.a() + b() * m.c(),
                a() * m.b() + b() * m.d(),
                c() * m.a() + d() * m.c(),
                c() * m.b() + d() * m.d()
            );
        } else {
            if constexpr (!isComplex && !std::is_integral_v<T>) {
                T p1 = (a() + d()) * (m.a() + m.d());
                T p2 = m.a() * (c() + d());
                T p3 = a() * (m.b() - m.d());
                T p4 = d() * (-m.a() + m.c());
                T p5 = m.d() * (a() + b());
                T p6 = (-a() + c()) * (m.a() + m.b());
                T p7 = (b() - d()) * (m.c() + m.d());
                
                return Matri2x2<T>(
                    p1 + p4 - p5 + p7, p3 + p5,
                    p2 + p4, p1 + p3 - p2 + p6
                );
            } else {
                return Matri2x2<T>(
                    a() * m.a() + b() * m.c(),
                    a() * m.b() + b() * m.d(),
                    c() * m.a() + d() * m.c(),
                    c() * m.b() + d() * m.d()
                );
            }
        }
    }

    [[nodiscard]] constexpr Matri2x2<T> operator * (const T& scalar) const noexcept {
        Matri2x2<T> result;
        result._data = _data * scalar;
        return result;
    }

    [[nodiscard]] constexpr Matri2x2<T> operator / (const T& scalar) const {
        if (std::abs(scalar) < _eps)
            throw std::runtime_error("Division by zero scalar!");
        [[assume(std::abs(scalar) >= _eps)]];
        return *this * (T(1) / scalar);
    }

    [[nodiscard]] friend constexpr Matri2x2<T> operator * (const T& scalar, const Matri2x2<T>& m) noexcept {
        return m * scalar;
    }

    [[nodiscard]] friend constexpr Matri2x2<T> operator / (const T& scalar, const Matri2x2<T>& m) {
        if (std::abs(m.determinant()) < m._eps)
            throw std::runtime_error("Matrix is singular!");
        return scalar * m.inverse();
    }

    constexpr Matri2x2<T>& operator+=(const Matri2x2<T>& m) noexcept {
        _data += m._data;
        return *this;
    }

    constexpr Matri2x2<T>& operator-=(const Matri2x2<T>& m) noexcept {
        _data -= m._data;
        return *this;
    }
    constexpr Matri2x2<T>& operator*=(const T& s) noexcept {
        _data *= s;
        return *this;
    }
    constexpr Matri2x2<T>& operator*=(const Matri2x2<T>& m) noexcept {
        *this = *this * m;
        return *this;
    }
    [[nodiscard]] constexpr T trace() const noexcept {
        return a() + d();
    }
    
    [[nodiscard]] constexpr T determinant() const noexcept {
        return a() * d() - b() * c();
    }
    
    [[nodiscard]] constexpr k_type norm() const noexcept {
        if constexpr (isComplex) {
            return std::sqrt(std::norm(a()) + std::norm(b()) + std::norm(c()) + std::norm(d()));
        } else {
            k_type sum = 0;
            for (size_t i = 0; i < 4; ++i)
                sum += _data[i] * _data[i];
            return std::sqrt(sum);
        }
    }

    void is_singular() const noexcept requires (!isComplex) {
        _singular = std::abs(this->determinant()) <= _eps;
    }

    void is_symmetric() const noexcept requires (!isComplex) {
        _symmetric = std::abs(b() - c()) <= _eps;
    }

    void is_nilpotent() const noexcept requires (!isComplex) {
        auto mat = (*this) * (*this);
        _nilpotent = std::abs(mat.a()) <= _eps && 
                      std::abs(mat.b()) <= _eps &&
                      std::abs(mat.c()) <= _eps && 
                      std::abs(mat.d()) <= _eps;
    }

    void is_idempotent() const noexcept requires (!isComplex) {
        auto mat = (*this) * (*this);
        _idempotent = std::abs(mat.a() - a()) <= _eps && 
                       std::abs(mat.b() - b()) <= _eps &&
                       std::abs(mat.c() - c()) <= _eps && 
                       std::abs(mat.d() - d()) <= _eps;
    }

    void is_involutory() const noexcept requires (!isComplex) {
        auto mat = (*this) * (*this);
        _involutory = std::abs(mat.a() - T(1)) <= _eps && 
                      std::abs(mat.b()) <= _eps &&
                      std::abs(mat.c()) <= _eps && 
                      std::abs(mat.d() - T(1)) <= _eps;
    }

    void is_orthogonal() const noexcept requires (!isComplex) {
        auto mt = this->transpose();
        auto mmt = mt * (*this);
        _orthogonal = std::abs(mmt.a() - T(1)) <= _eps && 
                       std::abs(mmt.b()) <= _eps &&
                       std::abs(mmt.c()) <= _eps && 
                       std::abs(mmt.d() - T(1)) <= _eps;
    }

    void is_upper_unipotent() const noexcept requires (!isComplex) {
        _upper_unipotent = std::abs(a() - T(1)) <= _eps && 
                           std::abs(d() - T(1)) <= _eps &&
                           std::abs(c()) <= _eps;
    }

    void is_diagonal_positive() const noexcept requires (!isComplex) {
        _diagonal_positive = std::abs(b()) <= _eps && 
                            std::abs(c()) <= _eps &&
                            a() > 0 && d() > 0;
    }

    void verifications() const noexcept requires (!isComplex) {
        is_singular();
        is_symmetric();
        is_nilpotent();
        is_idempotent();
        is_involutory();
        is_orthogonal();
        is_upper_unipotent();
        is_diagonal_positive();
    }

    [[nodiscard]] constexpr bool is_singular_acces() const noexcept requires (!isComplex) {
        return _singular.value_or(false);
    }

    [[nodiscard]] constexpr bool is_symmetric_acces() const noexcept requires (!isComplex) {
        return _symmetric.value_or(false);
    }

    [[nodiscard]] constexpr bool is_nilpotent_acces() const noexcept requires (!isComplex) {
        return _nilpotent.value_or(false);
    }

    [[nodiscard]] constexpr bool is_idempotent_acces() const noexcept requires (!isComplex) {
        return _idempotent.value_or(false);
    }

    [[nodiscard]] constexpr bool is_involutory_acces() const noexcept requires (!isComplex) {
        return _involutory.value_or(false);
    }

    [[nodiscard]] constexpr bool is_orthogonal_acces() const noexcept requires (!isComplex) {
        return _orthogonal.value_or(false);
    }

    [[nodiscard]] constexpr bool is_upper_unipotent_acces() const noexcept requires (!isComplex) {
        return _upper_unipotent.value_or(false);
    }

    [[nodiscard]] constexpr bool is_diagonal_positive_acces() const noexcept requires (!isComplex) {
        return _diagonal_positive.value_or(false);
    }

    [[nodiscard]] constexpr Matri2x2<T> transpose() const noexcept {
        return Matri2x2(a(), c(), b(), d());
    }

    [[nodiscard]] constexpr Matri2x2<T> adjugate() const noexcept {
        return Matri2x2(d(), -b(), -c(), a());
    }

    [[nodiscard]] Matri2x2<T> inverse() const {
        T det = determinant();
        
        k_type absDet;
        if constexpr (isComplex) {
            absDet = std::abs(det);
        } else {
            absDet = std::abs(det);
        }
        
        if (absDet <= _eps)
            throw std::runtime_error("Singular matrix !");
        T invDet = T(1) / det;
        
        return Matri2x2(d() * invDet, -b() * invDet, -c() * invDet, a() * invDet);
    }

    [[nodiscard]] auto Cholesky() const -> Matri2x2<T> requires (!isComplex) {
        if (!_symmetric.value_or(false))
            throw std::runtime_error("Unsymmetric matrix !");

        if (a() <= 0 || (a() * d() - b() * c()) <= 0)
            throw std::runtime_error("Not SPD matrix.");
        
        T l11 = std::sqrt(a());
        [[assume(l11 > T(0))]];
        T l21 = c() / l11;
        T l22 = std::sqrt(d() - l21 * l21);

        return Matri2x2<T>(l11, 0, l21, l22);
    }

    [[nodiscard]] auto LDLT() const -> std::tuple<Matri2x2<T>, Matri2x2<T>, Matri2x2<T>> requires (!isComplex) {
        if (!_symmetric.value_or(false))
            throw std::runtime_error("Unsymmetric matrix!");
        
        if (std::abs(a()) < _eps)
            throw std::runtime_error("Pivot too small or equal to zero!");
        
        T d1 = a();
        T l21 = b() / a();
        T d2 = d() - l21 * b();
        
        Matri2x2<T> L(T(1), T(0), l21, T(1));
        Matri2x2<T> D(d1, T(0), T(0), d2);
        Matri2x2<T> LT = L.transpose();
        
        return {L, D, LT};
    }

    [[nodiscard]] auto LU() const -> std::pair<Matri2x2<T>, Matri2x2<T>> {
        if (std::abs(a()) < _eps)
            throw std::runtime_error("Pivot too small or equal to zero!");
        [[assume(std::abs(a()) >= _eps)]];
        T l21 = c() / a();
        Matri2x2<T> L(T(1), T(0), l21, T(1));
        Matri2x2<T> U(a(), b(), T(0), d() - l21 * b());
        return {L, U};
    }

    [[nodiscard]] auto QR() const -> std::pair<Matri2x2<T>, Matri2x2<T>> requires (!isComplex) {
        T normeCol1 = std::hypot(a(), c());
        if (normeCol1 < _eps) {
            std::cerr << "Warning: First column is near zero, using default orthonormal basis." << std::endl;
            std::exit(EXIT_FAILURE);
        } else {
            [[assume(normeCol1 > _eps)]];
        }
        T q11 = (normeCol1 > _eps) ? a() / normeCol1 : T(0);
        T q21 = (normeCol1 > _eps) ? c() / normeCol1 : T(0);

        T dp = q11 * b() + q21 * d();
        T col2X = b() - dp * q11;
        T col2Y = d() - dp * q21;
        
        T normeCol2 = std::hypot(col2X, col2Y);
        T q12 = (normeCol2 > _eps) ? col2X / normeCol2 : T(0);
        T q22 = (normeCol2 > _eps) ? col2Y / normeCol2 : T(0);
        
        Matri2x2<T> Q(q11, q12, q21, q22);
        Matri2x2<T> R = Q.transpose() * (*this);

        return {Q, R};
    }

    [[nodiscard]] std::string characteristic_polynomial() const noexcept requires (!isComplex) {
        T det = determinant();
        T tr = trace();

        return std::format("P(x) = x²{}{}{}",
            tr > 0 ? std::format(" - {}x", tr) : (tr < 0 ? std::format(" + {}x", -tr) : ""),
            det > 0 ? std::format(" + {}", det) : (det < 0 ? std::format(" - {}", -det) : ""),
            (tr == 0 && det == 0) ? "0" : "");
    }

    [[nodiscard]] std::variant<
        std::pair<T, T>,
        std::pair<std::complex<T>, std::complex<T>>
    > eigenvalues() const noexcept requires (!isComplex) {
        const T tr = trace();
        const T det = determinant();
        const T disc = tr * tr - 4.0 * det;
        
        if (disc >= 0) {
            T sqrtDisc = std::sqrt(disc);
            T lambda1 = (tr - sqrtDisc) * 0.5;
            T lambda2 = (tr + sqrtDisc) * 0.5;
            return std::pair<T, T>(lambda1, lambda2);
        } else {
            T realPart = tr * 0.5;
            T imagPart = std::sqrt(-disc) * 0.5;
            return std::pair<std::complex<T>, std::complex<T>>(
                std::complex<T>(realPart, imagPart),
                std::complex<T>(realPart, -imagPart)
            );
        }
    }

    [[nodiscard]] std::variant<
        std::pair<std::array<T, 2>, std::array<T, 2>>,
        std::pair<std::array<std::complex<T>, 2>, std::array<std::complex<T>, 2>>
    > eigenvectors() const noexcept requires (!isComplex) {
        auto eigval = eigenvalues();
        
        if (std::holds_alternative<std::pair<T, T>>(eigval)) {
            auto [lambda1, lambda2] = std::get<0>(eigval);
            std::array<T, 2> v1, v2;
            
            if (std::abs(c()) > _eps) {
                v1[0] = d() - lambda1;
                v1[1] = -c();
            } else if (std::abs(b()) > _eps) {
                v1[0] = -b();
                v1[1] = a() - lambda1;
            } else {
                v1 = {1, 0};
            }
            
            T norm1 = std::hypot(v1[0], v1[1]);
            if (norm1 > _eps) {
                v1[0] /= norm1;
                v1[1] /= norm1;
            }
            
            if (std::abs(c()) > _eps) {
                v2[0] = d() - lambda2;
                v2[1] = -c();
            } else if (std::abs(b()) > _eps) {
                v2[0] = -b();
                v2[1] = a() - lambda2;
            } else {
                v2 = {0, 1};
            }
            
            T norm2 = std::hypot(v2[0], v2[1]);
            if (norm2 > _eps) {
                v2[0] /= norm2;
                v2[1] /= norm2;
            }
            
            return std::pair<std::array<T, 2>, std::array<T, 2>>(v1, v2);
        } else {
            auto [lambda1, lambda2] = std::get<1>(eigval);
            std::array<std::complex<T>, 2> v1;
            
            if (std::abs(c()) > _eps) {
                v1[0] = std::complex<T>(d(), 0) - lambda1;
                v1[1] = std::complex<T>(-c(), 0);
            } else if (std::abs(b()) > _eps) {
                v1[0] = std::complex<T>(-b(), 0);
                v1[1] = std::complex<T>(a(), 0) - lambda1;
            } else {
                v1 = {std::complex<T>(1, 0), std::complex<T>(0, 0)};
            }
            
            T norm = std::sqrt(std::norm(v1[0]) + std::norm(v1[1]));
            if (norm > _eps) {
                v1[0] /= norm;
                v1[1] /= norm;
            }
            
            std::array<std::complex<T>, 2> v2 = {std::conj(v1[0]), std::conj(v1[1])};
            
            return std::pair<std::array<std::complex<T>, 2>, std::array<std::complex<T>, 2>>(v1, v2);
        }
    }

    [[nodiscard]] auto diagonalization() const -> std::variant<
        std::tuple<Matri2x2<T>, Matri2x2<T>, Matri2x2<T>>,
        std::tuple<Matri2x2<complexType>, Matri2x2<complexType>, Matri2x2<complexType>>
    > requires (!isComplex) {
        auto eigval = eigenvalues();
        auto eigvec = eigenvectors();
        
        if (std::holds_alternative<std::pair<T, T>>(eigval)) {
            auto [lambda1, lambda2] = std::get<0>(eigval);
            auto [v1, v2] = std::get<0>(eigvec);
            
            T det_P = v1[0] * v2[1] - v1[1] * v2[0];
            if (std::abs(det_P) < _eps)
                throw std::runtime_error("-> Colinear eigenvectors!");
            
            Matri2x2<T> P(v1[0], v2[0], v1[1], v2[1]);
            Matri2x2<T> D(lambda1, 0, 0, lambda2);
            Matri2x2<T> invP = P.inverse();
            
            return std::tuple<Matri2x2<T>, Matri2x2<T>, Matri2x2<T>>{P, D, invP};
        } 
        else {
            auto [lambda1, lambda2] = std::get<1>(eigval);
            auto [v1, v2] = std::get<1>(eigvec);
            
            using cplx = complexType;
            
            Matri2x2<cplx> P(v1[0], v2[0], v1[1], v2[1]);
            Matri2x2<cplx> D(lambda1, cplx(0,0), cplx(0,0), lambda2);
            
            cplx detP = v1[0] * v2[1] - v1[1] * v2[0];
            if (std::abs(detP) < _eps)
               throw std::runtime_error("-> Colinear eigenvectors!");
            
            Matri2x2<cplx> invP = P.inverse();
            
            return std::tuple<Matri2x2<cplx>, Matri2x2<cplx>, Matri2x2<cplx>>{P, D, invP};
        }
    }

    [[nodiscard]] auto Schur() const -> std::pair<Matri2x2<T>, Matri2x2<T>> requires (!isComplex) {
        auto eigval = eigenvalues();
        
        if (std::holds_alternative<std::pair<std::complex<T>, std::complex<T>>>(eigval)) {
            auto [lambda1, lambda2] = std::get<1>(eigval);
            T alpha = lambda1.real();
            T beta = lambda1.imag();
            
            std::complex<T> a11 = std::complex<T>(a(), 0) - lambda1;
            std::complex<T> a12 = std::complex<T>(b(), 0);
            std::complex<T> a21 = std::complex<T>(c(), 0);
            std::complex<T> a22 = std::complex<T>(d(), 0) - lambda1;
            
            std::complex<T> v1, v2;
            if (std::abs(a12) > _eps) {
                v2 = std::complex<T>(1, 0);
                v1 = -a12 / a11;
            } else if (std::abs(a21) > _eps) {
                v1 = std::complex<T>(1, 0);
                v2 = -a21 / a22;
            } else {
                v1 = std::complex<T>(1, 0);
                v2 = std::complex<T>(0, 0);
            }
            
            T norm = std::sqrt(std::norm(v1) + std::norm(v2));
            if (norm > _eps) {
                v1 /= norm;
                v2 /= norm;
            }
            
            T q11 = v1.real(), q21 = v2.real();
            T q12 = v1.imag(), q22 = v2.imag();
            
            T normCol1 = std::hypot(q11, q21);
            T normCol2 = std::hypot(q12, q22);
            
            if (normCol1 > _eps) { q11 /= normCol1; q21 /= normCol1; }
            if (normCol2 > _eps) { q12 /= normCol2; q22 /= normCol2; }
            
            T dot = q11 * q12 + q21 * q22;
            if (std::abs(dot) > _eps) {
                q12 -= dot * q11;
                q22 -= dot * q21;
                T norm_new = std::hypot(q12, q22);
                if (norm_new > _eps) { q12 /= norm_new; q22 /= norm_new; }
            }
            
            Matri2x2<T> Q(q11, q12, q21, q22);
            Matri2x2<T> T_mat(alpha, beta, -beta, alpha);
            
            Matri2x2<T> test = Q * T_mat * Q.transpose();
            T error = std::abs(test(0,0) - (*this)(0,0)) + std::abs(test(0,1) - (*this)(0,1)) +
                      std::abs(test(1,0) - (*this)(1,0)) + std::abs(test(1,1) - (*this)(1,1));
            
            if (error > _eps)
                T_mat = Q.transpose() * (*this) * Q;
            
            return {Q, T_mat};
        }
        
        auto [lambda1, lambda2] = std::get<0>(eigval);
        
        bool swapNeeded = false;
        if (std::abs(lambda1) > std::abs(lambda2)) {
            std::swap(lambda1, lambda2);
            swapNeeded = true;
        }
        
        T a11 = a() - lambda1;
        T a12 = b();
        T a21 = c();
        T a22 = d() - lambda1;
        
        T vX, vY;
        if (std::abs(a12) > std::abs(a22) && std::abs(a12) > std::abs(a21)) {
            vY = 1;
            vX = -a12 / (a11 + _eps);
        } else if (std::abs(a21) > std::abs(a11) && std::abs(a21) > std::abs(a12)) {
            vX = 1;
            vY = -a21 / (a22 + _eps);
        } else if (std::abs(a11) > _eps) {
            vY = 1;
            vX = -a12 / a11;
        } else if (std::abs(a22) > _eps) {
            vX = 1;
            vY = -a21 / a22;
        } else {
            vX = 1;
            vY = 0;
        }
        
        T norm = std::hypot(vX, vY);
        if (norm > _eps) {
            vX /= norm;
            vY /= norm;
        }
        
        T uX = -vY;
        T uY = vX;
        
        Matri2x2<T> Q = swapNeeded ? Matri2x2<T>(uX, vX, uY, vY) : Matri2x2<T>(vX, uX, vY, uY);
        Matri2x2<T> T_mat = Q.transpose() * (*this) * Q;
        if (std::abs(T_mat(1,0)) < _eps)
            T_mat(1,0) = T(0);
        
        return {Q, T_mat};
    }

    [[nodiscard]] auto jordan_chevalley() const -> std::pair<Matri2x2<T>, Matri2x2<T>> requires (!isComplex) {
        const T tr = trace();
        const T det = determinant();
        const T disc = tr * tr - T(4) * det;

        if (disc >= 0) {
            // Real eigenvalues
            const T sqrtDisc = std::sqrt(disc);
            const T lambda1 = (tr - sqrtDisc) / T(2);
            const T lambda2 = (tr + sqrtDisc) / T(2);

            if (std::abs(lambda1 - lambda2) > _eps) {
                // Distinct eigenvalues => diagonalizable
                return {*this, Matri2x2<T>(T(0), T(0), T(0), T(0))};
            } else {
                // Double eigenvalue
                const T lambda = lambda1;
                const Matri2x2<T> d(lambda, T(0), T(0), lambda);
                const Matri2x2<T> n = (*this) - d;
                // Check if N² = 0
                const auto nn = n * n;
                const T nilpotence = std::abs(nn[0]) + std::abs(nn[1]) +
                                    std::abs(nn[2]) + std::abs(nn[3]);
                if (nilpotence < _eps && n.norm() > _eps) {
                    // Non-diagonalizable
                    return {d, n};
                } else {
                    // Diagonalizable (A = lambda I)
                    return {*this, Matri2x2<T>(T(0), T(0), T(0), T(0))};
                }
            }
        } else {
            return {*this, Matri2x2<T>(T(0), T(0), T(0), T(0))};
        }
    }
    [[nodiscard]] auto iwasawa() const -> std::tuple<Matri2x2<T>, Matri2x2<T>, Matri2x2<T>> requires (!isComplex) {
        if (std::abs(determinant()) < _eps)
            throw std::runtime_error("Singular matrix!");
        T normCol1 = std::hypot(a(), c());
        if (normCol1 < _eps)
            throw std::runtime_error("First column is zero!");
        [[assume(normCol1 >= _eps)]];

        T k11 = a() / normCol1;
        T k21 = c() / normCol1;        
        T dot = k11 * b() + k21 * d();
        T projX = dot * k11;
        T projY = dot * k21;
        T orthoX = b() - projX;
        T orthoY = d() - projY;
        
        T normOrtho = std::hypot(orthoX, orthoY);
        if (normOrtho < _eps) {
            k11 = 1.0; k21 = 0.0;
            orthoX = 0.0; orthoY = 1.0;
            normOrtho = 1.0;
        }
        
        T k12 = orthoX / normOrtho;
        T k22 = orthoY / normOrtho;
        
        // det(K) = -1 or 1
        Matri2x2<T> K(k11, k12, k21, k22);
        
        // If det(K) = -1, multiply the second column by -1
        T detK = K.determinant();
        if (detK < 0) {
            k12 = -k12;
            k22 = -k22;
            K = Matri2x2<T>(k11, k12, k21, k22);
        }
        
        // Étape 2: Compute A' = K^T * A
        Matri2x2<T> Aprime = K.transpose() * (*this);
        
        Aprime(1, 0) = T(0);
        
        // Étape 3: Extract A_diag and N de A'
        // A' = [[d1, x], [0, d2]] = [[d1, 0], [0, d2]] * [[1, n], [0, 1]]
        // where n = x / d1
        
        T d1 = Aprime(0, 0);
        T d2 = Aprime(1, 1);
        T x = Aprime(0, 1);
        
        if (std::abs(d1) < _eps || std::abs(d2) < _eps)
            throw std::runtime_error("Diagonal elements too small!");
        
        // Assurer d1, d2 > 0
        if (d1 < 0) {
            d1 = -d1;
            k11 = -k11;
            k21 = -k21;
            K = Matri2x2<T>(k11, k12, k21, k22);
        }
        
        if (d2 < 0) {
            d2 = -d2;
            k12 = -k12;
            k22 = -k22;
            K = Matri2x2<T>(k11, k12, k21, k22);
        }
        
        Matri2x2<T> diagA(d1, T(0), T(0), d2);
        T n = x / d1;
        Matri2x2<T> N(T(1), n, T(0), T(1));
        
        return {K, diagA, N};
    }

    [[nodiscard]] auto polar_decomposition() const -> std::pair<Matri2x2<T>, Matri2x2<T>> 
        requires (!isComplex) {
        if (!_symmetric.value_or(false))
            throw std::runtime_error("Unsymmetric matrix!");
        
        if (a() <= 0 || determinant() <= 0)
            throw std::runtime_error("Matrix must be positive definite!");
        
        T traceA = trace();
        T detA = determinant();    
        T sqrtDet = std::sqrt(detA);
        T s = std::sqrt(traceA + T(2) * sqrtDet);
        
        if (std::abs(s) < _eps)
            throw std::runtime_error("Polar decomposition: singular case!");
        
        Matri2x2<T> sqrtA;
        
        if (std::abs(s) > _eps) {
            T t = T(1) / s;
            sqrtA = (*this) + Matri2x2<T>(sqrtDet, T(0), T(0), sqrtDet);
            sqrtA *= t;
        } else {
            T lambda = std::sqrt(a());
            sqrtA = Matri2x2<T>(lambda, T(0), T(0), lambda);
        }
        
        Matri2x2<T> Q = Matri2x2<T>(T(1), T(0), T(0), T(1));
        
        return {Q, sqrtA};
    }

    [[nodiscard]] auto polar_decomposition_general() const
        -> std::pair<Matri2x2<T>, Matri2x2<T>> requires (!isComplex)
    {
        Matri2x2<T> mtm = this->transpose() * (*this);
        T a = mtm(0,0), b = mtm(0,1), d = mtm(1,1);

        T traceMtm = a + d;
        T detMtm = a * d - b * b;
        T disc = traceMtm * traceMtm - T(4) * detMtm;

        if (disc < 0)
            throw std::runtime_error("Polar decomposition: mtm has complex eigenvalues!");
        T sqrt_disc = std::sqrt(disc);
        T lambda1 = (traceMtm - sqrt_disc) / T(2);
        T lambda2 = (traceMtm + sqrt_disc) / T(2);

        if (lambda1 < 0 || lambda2 < 0)
            throw std::runtime_error("Polar decomposition: mtm not positive semidefinite!");

        T sqrt_lambda1 = std::sqrt(lambda1);
        T sqrt_lambda2 = std::sqrt(lambda2);

        Matri2x2<T> P;

        if (std::abs(b) < _eps) {
            // Diagonal case
            P = Matri2x2<T>(std::sqrt(a), T(0), T(0), std::sqrt(d));
        } else {
            // Eigenvector for lambda1
            T v1x, v1y;
            if (std::abs(a - lambda1) > _eps) {
                v1x = T(1);
                v1y = -(a - lambda1) / b;
            } else {
                v1y = T(1);
                v1x = -b / (d - lambda1);
            }
            T norm1 = std::hypot(v1x, v1y);
            v1x /= norm1;
            v1y /= norm1;

            // Eigenvector for lambda2
            T v2x = -v1y;
            T v2y =  v1x;

            // Rotation matrix R = [v1 v2]
            Matri2x2<T> R(v1x, v2x, v1y, v2y);
            Matri2x2<T> D_sqrt(sqrt_lambda1, T(0), T(0), sqrt_lambda2);

            // P = R * D_sqrt * R^T
            P = R * D_sqrt * R.transpose();
        }

        // U = M * P^{-1}
        Matri2x2<T> P_inv = P.inverse();
        Matri2x2<T> U = (*this) * P_inv;

        // Optional: force U to be exactly orthogonal (clean numerical noise)
        // U = (U + U.inverse().transpose()) / T(2);  // polar projection

        return {U, P};
    }
    [[nodiscard]] std::string display(int precision) const {
        if constexpr (isComplex) {
            auto format_complex = [precision](const T& c) -> std::string {
                k_type re = c.real();
                k_type im = c.imag();
                
                if (std::abs(im) < _eps) {
                    return std::format("{:.{}f}", re, precision);
                } else if (std::abs(re) < _eps) {
                    return std::format("{:.{}f}i", im, precision);
                } else if (im < 0) {
                    return std::format("{:.{}f}{:.{}f}i", re, precision, im, precision);
                } else {
                    return std::format("{:.{}f}+{:.{}f}i", re, precision, im, precision);
                }
            };
            
            return std::format("|{} {}|\n|{} {}|", 
                format_complex(a()), format_complex(b()),
                format_complex(c()), format_complex(d()));
        } else {
            return std::format("|{:.{}f} {:.{}f}|\n|{:.{}f} {:.{}f}|", 
                a(), precision, b(), precision, 
                c(), precision, d(), precision);
        }
    }

    friend std::ostream& operator<<(std::ostream& os, const Matri2x2<T>& m) {
        return os << m.display();
    }
};

// template<typename T>
// bool commutes(const Matri2x2<T>& A, const Matri2x2<T>& B) {
//     auto AB = A * B;
//     auto BA = B * A;
//     return (AB - BA).norm() < Matri2x2<T>::_eps;
// }

template<typename T> concept FloatingPrecision = std::same_as<T, float> || std::same_as<T, double> || std::same_as<T, long double>;

template<FloatingPrecision T> int get_precision() {
    if constexpr (std::same_as<T, float>) {
        int p = 6;
        std::println("Using single precision (float {})", p);
        return p;
    } else if constexpr (std::same_as<T, double>) {
        int p = 12;
        std::println("Using double precision (double {})", p);
        return p;
    } else if constexpr (std::same_as<T, long double>) {
        int p = 15;
        std::println("Using extended precision (long double {})", p);
        return p;
    } else {
        return 0;
    }
}

int main(int argc, char **argv) {
    using T = float; // float or double or long double
    T a = {}, b = {}, c = {}, d = {};

    int precision = get_precision<T>();
    if (precision == 0) {
        std::println(stderr, "Unsupported type for precision info.");
        std::exit(EXIT_FAILURE);
    }
        
    if (argc == 5) {
        auto parse = [](const char* str) -> T {
            char* end = nullptr;
            T coefs = std::strtof(str, &end);
            if (end == str || *end != '\0') {
                throw std::runtime_error(std::format("Invalid number: '{}'", str));
            }
            return coefs;
        };
        try {
            a = parse(argv[1]);
            b = parse(argv[2]);
            c = parse(argv[3]);
            d = parse(argv[4]);
        } catch (const std::runtime_error& e) {
            std::println(stderr, "Error parsing arguments: {}", e.what());
            std::exit(EXIT_FAILURE);
        }
    } else {
        std::random_device rnd;
        std::mt19937 gen(rnd());
        std::uniform_int_distribution<int> dist(-10, 10);
        a = dist(gen);
        b = dist(gen);
        c = dist(gen);
        d = dist(gen);
    }

    auto matrixPtr = std::make_unique<Matri2x2<T>>(a, b, c, d);
    Matri2x2<T>& matrix = *matrixPtr;
    
    if constexpr (!Matri2x2<T>::isComplex) {
        matrix.verifications();
    }

    std::println("\nMatrix:\n{}\n", matrix.display(precision));
    
    if constexpr (!Matri2x2<T>::isComplex) {
        std::println("Is singular ? {}", matrix.is_singular_acces() ? "Yes" : "No");
        std::println("Is symmetric ? {}", matrix.is_symmetric_acces() ? "Yes" : "No");
        std::println("Is nilpotent (A² = 0) ? {}", matrix.is_nilpotent_acces() ? "Yes" : "No");
        std::println("Is idempotent (A² = A) ? {}", matrix.is_idempotent_acces() ? "Yes" : "No");
        std::println("Is involutory (A² = I) ? {}", matrix.is_involutory_acces() ? "Yes" : "No");
        std::println("Is orthogonal (AA^T = I) ? {}\n", matrix.is_orthogonal_acces() ? "Yes" : "No");
    } else {
        std::unreachable();
    }
    
    std::println("Trace: {:.{}f}", matrix.trace(), precision);
    std::println("Determinant: {:.{}f}", matrix.determinant(), precision);
    std::println("Frobenius norm: {:.{}f}", matrix.norm(), precision);
    std::println();
    std::println("Adjugate:\n{}\n", matrix.adjugate().display(precision));

    std::string msgNonInv = "Singular matrix!\n)";
    if constexpr (!Matri2x2<T>::isComplex) {
        if (matrix.is_singular_acces())
            std::println("", msgNonInv);
    } else {
        if (std::abs(matrix.determinant()) <= Matri2x2<T>::_eps)
            std::println("", msgNonInv);
    }

    try {
        if (!matrix.is_singular_acces()) {
            std::println("Inverse:\n{}\n", matrix.inverse().display(precision));
            std::println("AA^(-1):\n{}\n", (matrix.inverse() * matrix).display(precision));
        }
    } catch (const std::runtime_error& e) {
        std::println(stderr, "Erreur: {}", e.what());
    }
        
    if constexpr (!Matri2x2<T>::isComplex) {
        if (matrix.is_symmetric_acces()) {
            std::println("Cholesky decomposition (symmetric matrix):");
            std::println("------------------------------------------");
            try {
                auto C = matrix.Cholesky();
                std::println("L:\n{}", C.display(precision));
                std::println("LL^T:\n{}\n", (C * C.transpose()).display(precision));
            } catch (const std::runtime_error& e) {
                std::println(stderr, "Cholesky: {}", e.what());
                std::println("Attempt with LDLT:");
                try {
                    auto [L, D, LT] = matrix.LDLT();
                    std::println("L:\n{}", L.display(precision));
                    std::println("D:\n{}", D.display(precision));
                    std::println("L^T:\n{}", LT.display(precision));
                    std::println("LDL^T:\n{}\n", (L * D * LT).display(precision));
                } catch (const std::runtime_error& e2) {
                    std::println(stderr, "LDLT: {}\n", e2.what());
                }
            }
        } else {
            std::println("LU decomposition (unsymmetric matrix):");
            std::println("--------------------------------------");
            try {
                auto [L, U] = matrix.LU();
                std::println("L:\n{}", L.display(precision));
                std::println("U:\n{}", U.display(precision));
                std::println("LU:\n{}\n", (L * U).display(precision));
            } catch (const std::runtime_error& e) {
                std::println(stderr, "LU: {}\n", e.what());
            }
        }

        std::println("QR decomposition:");
        std::println("-----------------");
        try {
            auto [Q, R] = matrix.QR();
            std::println("Q:\n{}", Q.display(precision));
            std::println("R:\n{}", R.display(precision));
            std::println("QR:\n{}\n", (Q * R).display(precision));
        } catch (const std::runtime_error& e) {
            std::println(stderr, "QR: {}\n", e.what());
        }

        std::println("Characteristic polynomial:");
        std::println("--------------------------");
        std::println("{}", matrix.characteristic_polynomial());

        auto valp = matrix.eigenvalues();
        if (std::holds_alternative<std::pair<T, T>>(valp)) {
            auto [l1, l2] = std::get<0>(valp);
            std::println("Real eigenvalues (roots): {:.{}f}, {:.{}f}\n", l1, precision, l2, precision);
        } else {
            auto [l1, l2] = std::get<1>(valp);
            std::println("Complex eigenvalues: {}, {}\n", l1, l2);
        }

        auto vecp = matrix.eigenvectors();
        if (std::holds_alternative<std::pair<std::array<T, 2>, std::array<T, 2>>>(vecp)) {
            auto [v1, v2] = std::get<0>(vecp);
            std::println("Real eigenvectors:");
            std::println("  v1 = [{:.{}f}, {:.{}f}]", v1[0], precision, v1[1], precision);
            std::println("  v2 = [{:.{}f}, {:.{}f}]\n", v2[0], precision, v2[1], precision);
        } else {
            auto [v1, v2] = std::get<1>(vecp);
            std::println("Complex eigenvectors:");
            std::println("  v1 = [{}, {}]", v1[0], v1[1]);
            std::println("  v2 = [{}, {}]\n", v2[0], v2[1]);
        }

        std::println("\nDiagonalization:");
        std::println("----------------");
        try {
            auto result = matrix.diagonalization();
            
            if (std::holds_alternative<std::tuple<Matri2x2<T>, Matri2x2<T>, Matri2x2<T>>>(result)) {
                auto [P, D, invP] = std::get<0>(result);
                std::println("Passage matrix P:\n{}\n", P.display(precision));
                std::println("Diagonal matrix D:\n{}\n", D.display(precision));
                std::println("P^(-1):\n{}\n", invP.display(precision));
                std::println("A = PDP^(-1):\n{}\n", (P * D * invP).display(precision));
            } else {
                auto [P, D, invP] = std::get<1>(result);
                std::println("Passage matrix P (complex):\n{}\n", P.display(precision));
                std::println("Diagonal matrix D (complex):\n{}\n", D.display(precision));
                std::println("P^(-1) (complex):\n{}\n", invP.display(precision));
                std::println("A = PDP^(-1):\n{}", (P * D * invP).display(precision));
            }
        } catch (const std::runtime_error& e) {
            std::println(stderr, "{}", e.what());
        }

        std::println("\nSchur decomposition (A = Q*T*Q^T):");
        std::println("----------------------------------");
        try {
            auto [Q, T_mat] = matrix.Schur();
            std::println("Q (orthogonal):\n{}\n", Q.display(precision));
            std::println("T (upper triangular):\n{}\n", T_mat.display(precision));
            std::println("Checking QTQ^T:\n{}\n", (Q * T_mat * Q.transpose()).display(precision));
        } catch (const std::runtime_error& e) {
            std::println(stderr, "Schur: {}", e.what());
        }
        
        std::println("\nJordan-Chevalley Decomposition (Dunford) (A = D + N, DN = ND):");
        std::println("--------------------------------------------------------------");
        try {
            auto [D, N] = matrix.jordan_chevalley();
            std::println("D:\n{}", D.display(precision));
            std::println("N:\n{}", N.display(precision));
            
            std::println("\nVerification:");
            std::println("A = D + N:\n{}", (D + N).display(precision));
            std::println("N² = 0:\n{}", (N * N).display(precision));
            
            auto DN = D * N;
            auto ND = N * D;
            std::println("DN - ND:\n{}", (DN - ND).display(precision));            
        } catch (const std::exception& e) {
            std::println(stderr, "Jordan-Chevalley decomposition error: {}", e.what());
        }
    } else {
        std::unreachable();
    }

    std::println("\nIwasawa Decomposition (A = KDN)");
    std::println("-------------------------------");
    try {
        auto [K, D, N] = matrix.iwasawa();
        
        std::println("K:");
        std::println("{}", K.display(precision));
        std::println("det(K) = {:.{}f}", K.determinant(), precision);
        
        std::println("\nD:");
        std::println("{}", D.display(precision));
        
        std::println("\nN:");
        std::println("{}", N.display(precision));
        
        std::println("\nVerification A = KDN:");
        Matri2x2<T> KDN = K * D * N;
        std::println("K * D * N:\n{}", KDN.display(precision));
        
        T error = (matrix - KDN).norm();
        std::println("||A - KDN|| = {:.{}e}", error, precision);
        
        std::println("\nProperties verification:");
        std::println("K orthogonal? {}", K.is_orthogonal_acces() ? "Yes" : "No");
        D.verifications();
        N.verifications();
        std::println("D diagonal positive? {}", D.is_diagonal_positive_acces() ? "Yes" : "No");
        std::println("N upper unipotent? {}", N.is_upper_unipotent_acces() ? "Yes" : "No");
    } catch (const std::runtime_error& e) {
        std::println(stderr, "Iwasawa: {}", e.what());
    }

    std::println("\nPolar Decomposition (A = UP):");
    std::println("-----------------------------");

    if (matrix.is_symmetric_acces()) {
        try {
            // Version for symmetric positive definite matrices
            auto [U, P] = matrix.polar_decomposition();
            std::println("Orthogonal part U:\n{}", U.display(precision));
            std::println("Symmetric positive part P:\n{}", P.display(precision));
            
            std::println("\nVerification A = UP:");
            std::println("UP:\n{}", (U * P).display(precision));
            
            std::println("\nVerification P² = A:");
            std::println("P²:\n{}", (P * P).display(precision));
            
        } catch (const std::runtime_error& e) {
            std::println(stderr, "Polar decomposition (symmetric): {}", e.what());
            
            std::println("\n=== General Polar Decomposition ===");
            try {
                auto [U, P] = matrix.polar_decomposition_general();
                std::println("Orthogonal part U:\n{}", U.display(precision));
                std::println("Symmetric positive part P:\n{}", P.display(precision));
                
                std::println("\nVerification A = UP:");
                std::println("UP:\n{}", (U * P).display(precision));
                std::println("U orthogonal? {}", 
                    (U.transpose() * U == Matri2x2<T>(T(1), T(0), T(0), T(1))) ? "Yes" : "No");
                std::println("||A - UP|| = {:.{}e}", (matrix - U * P).norm(), precision);
                
            } catch (const std::runtime_error& e2) {
                std::println(stderr, "General polar: {}", e2.what());
            }
        }
    } else {
        std::println("Matrix not symmetric - using general polar decomposition:");
        try {
            auto [U, P] = matrix.polar_decomposition_general();
            std::println("Orthogonal part U:\n{}", U.display(precision));
            std::println("Symmetric positive part P:\n{}", P.display(precision));
            
            std::println("\nVerification A = UP:");
            std::println("UP:\n{}", (U * P).display(precision));
            std::println("||A - UP|| = {:.{}e}", (matrix - U * P).norm(), precision);
            std::println("U orthogonal? {}", 
                (U.transpose() * U == Matri2x2<T>(T(1), T(0), T(0), T(1))) ? "Yes" : "No");
            
        } catch (const std::runtime_error& e) {
            std::println(stderr, "Polar decomposition: {}", e.what());
        }
    }

    std::fflush(stdout);
    std::fflush(stderr);

    matrixPtr = nullptr;

    return EXIT_SUCCESS;
}
\end{lstlisting}

\end{document}
